\documentclass[a4j,11pt,dvipdfmx]{jsarticle}
\usepackage{float}
\usepackage{array}
\usepackage{titlesec}
\usepackage[dvipdfmx]{graphicx}
\usepackage{graphicx}
\usepackage{comment}
\usepackage{multirow}
\usepackage{array}
\usepackage{longtable}

\titleformat*{\section}{\bfseries}
\titleformat*{\subsection}{\bfseries}

\makeatletter
  \renewcommand{\section}{%
    \@startsection{section}{1}{\z@}%
    {0.4\Cvs}{0.1\Cvs}%
    {\normalfont\headfont\raggedright}}
\makeatother

\makeatletter
  % sectionの下マージンを小さく
  \renewcommand{\subsection}{%
    \@startsection{subsection}{1}{\z@}%
    {0.1\Cvs}{0.1\Cvs}%
    {\normalfont\headfont\raggedright}}
\makeatother

\renewcommand{\thesubsection}{\thesection-\arabic{subsection}.}


%---------------------------------------------------
% ページの設定
%---------------------------------------------------
\setlength{\textwidth}{160truemm}
\setlength{\textheight}{250truemm}
\setlength{\topmargin}{-4.5truemm}
\setlength{\oddsidemargin}{0.5truemm}
\pagestyle{empty}
\setlength{\headheight}{0truemm}
\setlength{\parindent}{1zw}
\newcolumntype{b}{!{\vrule width 1pt}}
\newcommand{\bhline}{\noalign{\hrule height 1pt}}


\begin{document}
提出日: 2026年4月29日
\begin{center}
%\huge{実験計画書}
\huge{第3回ハッカソン議事録}
\end{center}

\begin{table}[htbp]
    \begin{tabular}{llrr}
    \textgt{開催日時}: & 2026年4月29日 12:00〜16:00& &\\
    \textgt{開催場所}: &津田沼校舎 7号館3F 732講義室& &\\

     \textgt{班名}:& 23班 & & \\
     \textgt{メンバー}:& 24G1038 片岡聖  (議事録作成者)& & \\%k
                      & 24G1131 地曳賢人 & & \\%z
                    & 24G1131 御代川稜 & & \\%m
                    & 25G1021 梅澤颯太 & & \\%u
                    & 25G1053 齋藤翔太 & & \\%s
                    & 25G1065 塩澤匠生 & & \\%ta
        \textgt{欠席者}: & 未定 & &  \\
        \textgt{目的}: &FBS, PBS, WBS, MOP, V\&V, TPM, \\
        &役割分担のチーム内で確認，修正，確定.  & & \\
        &WBS と役割分担をもとに， スケジュールを作成. \\
        &基本設計と詳細設計について設計に対するレビュー.\\

    \end{tabular}
\end{table}

%%%%%%%%%%%%%%%%%%%%%%%%%%%%%%%%%%%%%%%%%%%%%%%%%
%%%%%%%%%%%%%%%%%%%%%%%%%%%%%%%%%%%%%%%%%%%%%%%%%
%%%%%%%%%%%%%%%%%%%%%%%%%%%%%%%%%%%%%%%%%%%%%%%%%

\section{議題一覧}
・FBS, PBS, WBS,役割分担のチーム内で確認，修正，確定
・事前課題のレビュー
\section{討議内容}

\subsection{FBS,PBS,WBS,役割分担の確認}

塩澤:事前に考えてきたWBS図を作成し，gitにプッシュしておく

御代川:塩澤に負担がかかりすぎている気がする

地曳:楽器担当とリズム担当を分ける

塩澤:分析，解析を齊藤，prosessingを倍音など土台作成を梅澤に振り直すのはどうか

御代川:作り直したWBS図を元にスケジュール組み直す

御代川:それぞれの工程がどれぐらいかかるかの確認(分かる範囲で)

片岡:製造自体のフェーズのかかる時間を短く設定して，遅れた分を予備日にまわす

御代川，塩澤:最低でもテストはしたいため，早めの完成を予定する

塩澤:プレゼンまでに最低でも基本設計，詳細設計ちょっとまで完成させたほうがよい

塩澤:テストフェーズで評価指標との比較をするため時間を取りたい

塩澤:設定3週間，製造単体結合2週間，テストフェーズ2週間で行いたい


\subsection{基本設計と詳細設計について設計に対するレビュー}

\textgt{片岡聖の事前課題のレビュー}

平均評価スコア1/10

\textgt{地曳賢人の事前課題のレビュー}

平均評価スコア6/10

\textgt{御代川稜の事前課題のレビュー}

平均評価スコア3/10


\textgt{梅澤颯太の事前課題のレビュー}

平均評価スコア3.5/10

\textgt{齋藤翔太の事前課題のレビュー}

平均評価スコア2.5/10

\textgt{塩澤匠生の事前課題のレビュー}

平均評価スコア9.5/10










\subsection{個人毎の次回までに実施する内容}

\textgt{片岡聖，地曳賢人，御代川稜}

・MOP, V\&V, TPMのの検討

・ガントチャートの作成

・担当箇所の基本設計・詳細設計の作成

\textgt{梅澤颯太}

・音色合成基本設計

・担当箇所の基本設計・詳細設計の作成

\textgt{齋藤翔太}

・楽譜データ形成基本設計

・担当箇所の基本設計・詳細設計の作成

\textgt{塩澤匠生}

・システムアーキテクチャ設計

・通信プロトコル基本設計

・担当箇所の基本設計・詳細設計の作成

\section{決定事項}

最終的な役割分担，またそれぞれの工程でどれほど時間がかかるかの表は以下の表\ref{yakuwari}，\ref{yakuwari_2}
に記載

\begin{table}[htbp]
\centering
\caption{新たに作成した作業分担表}
\renewcommand{\arraystretch}{0.9}
\label{yakuwari}
\begin{tabular}{|c|c|c|p{5cm}|c|c|}
\hline
時間的フェーズ & 番号 & 要素成果物 & 活動タスク & 担当者 & 必要時間 \\ \hline

% 設計フェーズ（合計14行）
\multirow{14}{*}{\shortstack{100 設計フェーズ \\ (3週間)}}
& 110 & 基本設計 & 111 FBSの最終化 & 全員 & - \\ \cline{4-6}
& & & 112 PBSの最終化 & 全員 & - \\ \cline{4-6}
& & & 112 MOE/MOP/TPMの最終化 & 全員 & - \\ \cline{4-6}
& & & 111 システムアーキテクチャ設計 & 塩澤 & 1週間 \\ \cline{4-6}
& & & 112 通信プロトコル基本設計 & 塩澤 & 1週間 \\ \cline{4-6}
& & & 113 楽譜データ形式基本設計 & 齊藤 & 2週間 \\ \cline{4-6}
& & & 114 音色合成基本設計 & 梅澤 & 2週間 \\ \cline{2-6}
& 120 & 詳細設計 & 121 共通層API詳細 & 塩澤 & - \\ \cline{4-6}
& & & 122 指揮者の詳細設計 & 塩澤 & - \\ \cline{4-6}
& & & 123 楽器の詳細設計 & 塩澤 & - \\ \cline{4-6}
& & & 124 Arduinoから届く司令で鳴らす仕組み & 梅澤 & - \\ \cline{4-6}
& & & 125 4種類の楽器パートを鳴らす仕組み & 梅澤 & 2週間 \\ \cline{4-6}
& & & 126 楽器のような音色の作成 & 梅澤 & 2週間 \\ \cline{4-6}
& & & 127 楽器音を周波数分解し126にフィードバック & 齊藤 & 2週間 \\ \cline{4-6}
& & & 128 楽器パートの振り分け & 斎藤 & 3週間 \\ \cline{4-6}
& & & 129 テンポ，音の高さ，リズムの形の決定 & 斎藤 & 3週間 \\ \cline{4-6}
& & & 130 ArduinoからPCへ送る音データの中身 & 斎藤 & 3週間 \\ \hline

% 製造フェーズ（合計19行）
\multirow{19}{*}{\shortstack{200 製造フェーズ \\ (2週間)}}
& 210 & 共通層実装 & 211 IModule/ModuleTimer実装 & 塩澤 & - \\ \cline{4-6}
& & & 212 SystemData/ProjectConfig テンプレ整備 & 塩澤 & - \\ \cline{4-6}
& & & 213 OrcNetModule & 塩澤 & - \\ \cline{2-6}
& 220 & 指揮者 & 221 IMUドライバモジュール & 塩澤 & - \\ \cline{4-6}
& & & 222 信号前処理モジュール & 塩澤 & - \\ \cline{4-6}
& & & 223 拍検出モジュール & 塩澤 & - \\ \cline{4-6}
& & & 224 テンポ推定モジュール & 塩澤 & - \\ \cline{4-6}
& & & 225 コマンド送出モジュール & 塩澤 & - \\ \cline{2-6}
& 230 & 楽器 & 231 CTRL/BEAT受信モジュール & 塩澤 & - \\ \cline{4-6}
& & & 232 楽譜データ保持・進行ロジック & 塩澤 & - \\ \cline{4-6}
& & & 233 NOTE送出モジュール & 塩澤 & - \\ \cline{4-6}
& & & 234 パート別差分適用 & 塩澤 & - \\ \cline{2-6}
& 240 & 楽譜準備 & 241 課題曲4パートの選定 & 齊藤 & 2週間 \\ \cline{4-6}
& & & 242 4パート楽譜のヘッダ配列化 & 齊藤 & 2週間 \\ \cline{4-6}
& & & 243 楽譜データのROM埋め込み確認 & 齊藤 & 2週間 \\ \cline{2-6}
& 250 & Processing実装 & 251 NOTE受信モジュール & 梅澤 & 2週間 \\ \cline{4-6}
& & & 252 金管音色合成エンジン & 梅澤 & 2週間 \\ \cline{4-6}
& & & 253 4パート同時発音管理 & 梅澤 & 2週間 \\ \cline{4-6}
& & & 254 UI/モード表示 & 梅澤 & 2週間 \\ \hline

\end{tabular}
\end{table}

\begin{table}[htbp]
\centering
\caption{新たに作成した作業分担表の続き}
\renewcommand{\arraystretch}{0.8}
\label{yakuwari_2}
\begin{tabular}{|c|c|c|p{5cm}|c|c|}
\hline
時間的フェーズ & 番号 & 要素成果物 & 活動タスク & 担当者 & 必要時間 \\ \hline

% テストフェーズ
\multirow{15}{*}{\shortstack{300 テストフェーズ \\ (2週間)}}
& 310 & 単体テスト & 311 共通層単体 & 塩澤 & - \\ \cline{4-6}
& & & 312 拍検出ゴールデンテスト & 塩澤 & - \\ \cline{4-6}
& & & 313 楽譜進行状態機械テスト & 塩澤 & - \\ \cline{4-6}
& & & 314 音色生成単体確認 & 梅澤 & 2週間 \\ \cline{2-6}
& 320 & 結合テスト & 321 指揮単体デモ & 塩澤 & - \\ \cline{4-6}
& & & 322 指揮者と楽器の結合 & 塩澤 & - \\ \cline{4-6}
& & & 323 ４ノード同期誤差計測 & 全員 & - \\ \cline{4-6}
& & & 324 全ノードとProcessing結合 & 全員 & - \\ \cline{2-6}
& 330 & システムテスト & 331 拍検出精度 & 塩澤 & - \\ \cline{4-6}
& & & 332 通信遅延 & 塩澤 & - \\ \cline{4-6}
& & & 333 動機誤差 & 全員 & - \\ \cline{4-6}
& & & 334 テンポ追従 & 塩澤 & - \\ \cline{4-6}
& & & 335 パケロス耐性 & 塩澤 & - \\ \cline{4-6}
& & & 336 CPU負荷 & 塩澤 & - \\ \cline{4-6}
& & & 337 起動時間 & 塩澤 & - \\ \hline

\multirow{3}{*}{400 ストレッチフェーズ}
& 410 & 強弱推定 & 411 velocity推定アルゴ実装 & 塩澤 & - \\ \cline{4-6}
& & & 412 CTRLパケットへのvelocity乗せこみ & 塩澤 & - \\ \cline{4-6}
& & & 413 強弱制御テスト & 塩澤 & - \\ \hline

\multirow{2}{*}{500 運用・発表}
& 510 & 発表 & 511 発表準備・デモ調整 & 全員 & - \\ \cline{2-6}
& 520 & 報告書 & 521 成果報告書(個人) & 全員 & - \\ \hline

\end{tabular}
\end{table}

\section{継続審議事項・保留事項}

・MOP, V\&V, TPMの確認，修正，確定
・ガントチャートは作成途中


\section{次回の会議の日時・場所・予定議題}

\begin{table}[htbp]
    \begin{tabular}{llrr}
    \textgt{開催日時}: & 2026年5月13日12:00〜16:00 & &\\
    \textgt{開催場所}: &津田沼校舎 7号館3F 732講義室& &\\

        \textgt{予定議題}: &チーム全体のプロジェクトの計画書とシステムの設計書，\\
                &及び，プレゼンテーションについても議論&\\
    \end{tabular}
\end{table}

\section{その他，特記事項}
%未定
特になし

\end{document}
